\subsection{Funtores de hipertor global para complejos de módulos
cuasicoherentes y sucesiones espectrales de Künneth: caso de base afín}
\label{subsection:III.6.6}

\begin{env}[6.6.1]
\label{III.6.6.1}
Consideremos un esquema afín $S=\operatorname{Spec}(A)$ y dos $S$-esquemas
cuasicompactos $X^{(i)}$ ($i=1,2$); sea $\mathscr P_\bullet^{(i)}$ un complejo
de $\mathscr O_{X^{(i)}}$-módulos cuasicoherentes, acotado inferiormente,
cuyo diferencial es de grado $-1$ ($i=1,2$). Consideremos, por otra parte, un
cubrimiento finito
$\mathfrak U^{(i)}=(U_\alpha^{(i)})$ de $X^{(i)}$ por abiertos afines; sea
\[
  X=X^{(1)}\times_S X^{(2)},
\]
que es un $S$-esquema cuasicompacto (I, 5.5.1 y 6.6.4), y sea
$\mathfrak U=\mathfrak U^{(1)}\times_S\mathfrak U^{(2)}$ el cubrimiento de
$X$ formado por los abiertos afines
$U_\alpha^{(1)}\times_S U_\beta^{(2)}$. Para todo par de enteros
$p\leqslant0$, $q\in\mathbf Z$, el grupo de las $(-p)$-cocadenas alternadas
\[
  C^{-p}(\mathfrak U^{(i)},\mathscr P_q^{(i)})
\]
del cubrimiento $\mathfrak U^{(i)}$ con coeficientes en el haz
$\mathscr P_q^{(i)}$ (G, II, 5.1) es un $A$-módulo; para $p>0$, pondremos
$C^p(\mathfrak U^{(i)},\mathscr P_q^{(i)})=0$; queda así definido un
bicomplejo
\[
  C^\bullet(\mathfrak U^{(i)},\mathscr P_\bullet^{(i)})
\]
de $A$-módulos, cuyos dos diferenciales son de grado $-1$. De las
definiciones ($0_{\mathrm I}$, 12.1.2) resulta que el $A$-módulo de homología
\[
  \mathrm H_n\bigl(C^\bullet
  (\mathfrak U^{(i)},\mathscr P_\bullet^{(i)})\bigr)
\]
de este bicomplejo (considerado, como de costumbre, como un complejo simple
para el grado total) no es sino el $A$-módulo de hipercohomología
\[
  \mathbf H^{-n}
  (\mathfrak U^{(i)},\mathscr P^{(i)\bullet}),
\]
donde $\mathscr P^{(i)\bullet}$ es el complejo con diferencial de grado $+1$
obtenido tomando $\mathscr P_{-q}^{(i)}$ como componente de grado $q$; por
abuso de notación, lo escribiremos
$\mathbf H^{-n}(\mathfrak U^{(i)},\mathscr P_\bullet^{(i)})$. Resulta entonces
de (6.2.2) que este $A$-módulo es canónicamente isomorfo al $A$-módulo de
hipercohomología
$\mathbf H^{-n}(X^{(i)},\mathscr P^{(i)\bullet})$, que escribiremos asimismo
$\mathbf H^{-n}(X^{(i)},\mathscr P_\bullet^{(i)})$; por tanto, no depende del
cubrimiento finito $\mathfrak U^{(i)}$ elegido.
\end{env}

\begin{env}[6.6.2]
\label{III.6.6.2}
A los dos bicomplejos de $A$-módulos
\[
  L_{\bullet\bullet}^{(i)}
  =C^\bullet(\mathfrak U^{(i)},\mathscr P_\bullet^{(i)})
  \qquad(i=1,2)
\]
y al bifuntor covariante que asocia a estos dos bicomplejos su producto
tensorial $L_{\bullet\bullet}^{(1)}\otimes_A L_{\bullet\bullet}^{(2)}$ les
aplicaremos la teoría general de la hiperhomología de funtores respecto de
bicomplejos ($0_{\mathrm I}$, 11.7.4). Como las cocadenas consideradas son
alternadas y los cubrimientos $\mathfrak U^{(i)}$ son finitos, observemos que
los módulos $C^{-p}(\mathfrak U^{(i)},\mathscr P_q^{(i)})$ solo son $\ne0$
para un número finito (independiente de $q$) de valores de $p$ y, en
particular, las dos graduaciones de cada $L_{\bullet\bullet}^{(i)}$ están
acotadas inferiormente. Denotaremos por
\[
  \mathbf{Tor}_n^A(L_{\bullet\bullet}^{(1)},L_{\bullet\bullet}^{(2)})
\]
o por
\[
  \mathbf{Tor}_n^S
  (\mathfrak U^{(1)},\mathfrak U^{(2)};
   \mathscr P_\bullet^{(1)},\mathscr P_\bullet^{(2)})
\]
el $n$-ésimo módulo de hiperhomología de
$L_{\bullet\bullet}^{(1)}\otimes_A L_{\bullet\bullet}^{(2)}$, al que
llamaremos el \emph{hipertor} de índice $n$ de
$\mathscr P_\bullet^{(1)}$ y $\mathscr P_\bullet^{(2)}$, relativo a los
cubrimientos $\mathfrak U^{(1)}$ y $\mathfrak U^{(2)}$. Cuando
$\mathscr P_\bullet^{(1)}$ y $\mathscr P_\bullet^{(2)}$ están reducidos a sus
términos de grado 0, $\mathscr F^{(1)}$ y $\mathscr F^{(2)}$, escribimos
\[
  \mathbf{Tor}_n^S
  (\mathfrak U^{(1)},\mathfrak U^{(2)};
   \mathscr F^{(1)},\mathscr F^{(2)})
\]
para su hipertor. Conforme a las convenciones generales, denotaremos por
\[
  \mathbf{Tor}_n^S
  (\mathfrak U^{(1)},\mathfrak U^{(2)};
   \mathscr P_\bullet^{(1)},\mathscr P_\bullet^{(2)})
\]
el bicomplejo cuya componente de índices $(j,k)$ es
\[
  \mathbf{Tor}_n^S
  (\mathfrak U^{(1)},\mathfrak U^{(2)};
   \mathscr P_j^{(1)},\mathscr P_k^{(2)}).
\]

Como $L_{\bullet\bullet}^{(i)}$ es un funtor exacto respecto de
$\mathscr P_\bullet^{(i)}$, puesto que las intersecciones de los miembros
\oldpage[III]{154}
de $\mathfrak U^{(i)}$ son afines (I, 5.5.6 y 1.3.11),
$\mathbf{Tor}_\bullet^S
(\mathfrak U^{(1)},\mathfrak U^{(2)};
 \mathscr P_\bullet^{(1)},\mathscr P_\bullet^{(2)})$ es un
bi-$\partial$-funtor covariante en $\mathscr P_\bullet^{(1)}$ y
$\mathscr P_\bullet^{(2)}$, con valores en la categoría de los $A$-módulos
($0_{\mathrm I}$, 11.7.3). Además, se sabe ($0_{\mathrm I}$, 11.7.2) que este
bifuntor es el término límite común de seis funtores espectrales birregulares,
que denotaremos por
\[
  {}^{(t)}\mathscr E^S
  (\mathfrak U^{(1)},\mathfrak U^{(2)};
   \mathscr P_\bullet^{(1)},\mathscr P_\bullet^{(2)})
\]
o por
\[
  {}^{(t)}\mathscr E
  (\mathfrak U^{(1)},\mathfrak U^{(2)};
   \mathscr P_\bullet^{(1)},\mathscr P_\bullet^{(2)}),
\]
donde $t$ ha de sustituirse por una de las letras $a$, $b$, $a'$, $b'$, $c$,
$d$, y cuyos términos $E_2$ son los siguientes:
\[
\begin{aligned}
  {}^{(a)}\mathscr E_{pq}^2
  &=\bigoplus_{q_1+q_2=q}
    \mathbf{Tor}_p^A
    \bigl(\mathrm H_{q_1}(L_{\bullet\bullet}^{(1)}),
          \mathrm H_{q_2}(L_{\bullet\bullet}^{(2)})\bigr),\\
  {}^{(b)}\mathscr E_{pq}^2
  &=\mathrm H_p\bigl(
    \mathbf{Tor}_q^{A,\mathrm{II}}
    (L_{\bullet\bullet}^{(1)},L_{\bullet\bullet}^{(2)})\bigr),\\
  {}^{(a')}\mathscr E_{pq}^2
  &=\bigoplus_{q_1+q_2=q}
    \mathbf{Tor}_p^A
    \bigl(\mathrm H_{q_1}^{\mathrm I}(L_{\bullet\bullet}^{(1)}),
          \mathrm H_{q_2}^{\mathrm I}(L_{\bullet\bullet}^{(2)})\bigr),\\
  {}^{(b')}\mathscr E_{pq}^2
  &=\mathrm H_p\bigl(
    \mathbf{Tor}_q^A
    (L_{\bullet\bullet}^{(1)},L_{\bullet\bullet}^{(2)})\bigr),\\
  {}^{(c)}\mathscr E_{pq}^2
  &=\bigoplus_{q_1+q_2=q}
    \mathbf{Tor}_p^A
    \bigl(\mathrm H_{q_1}^{\mathrm{II}}(L_{\bullet\bullet}^{(1)}),
          \mathrm H_{q_2}^{\mathrm{II}}(L_{\bullet\bullet}^{(2)})\bigr),\\
  {}^{(d)}\mathscr E_{pq}^2
  &=\mathrm H_p\bigl(
    \mathbf{Tor}_q^{A,\mathrm I}
    (L_{\bullet\bullet}^{(1)},L_{\bullet\bullet}^{(2)})\bigr),
\end{aligned}
\]
donde las notaciones concuerdan con las de la teoría general de la
hiperhomología. A continuación explicitaremos con mayor detalle estos términos
iniciales.
\end{env}

\begin{env}[6.6.3]
\label{III.6.6.3}
\emph{Sucesiones espectrales $(a)$ y $(a')$.} Hemos visto en (6.6.1) que el
módulo de homología $\mathrm H_n(L_{\bullet\bullet}^{(i)})$ del bicomplejo
$L_{\bullet\bullet}^{(i)}$ era igual a
$\mathbf H^{-n}(X^{(i)},\mathscr P_\bullet^{(i)})$; por tanto,
\[
  {}^{(a)}\mathscr E_{pq}^2
  =\bigoplus_{q_1+q_2=q}
   \mathbf{Tor}_p^A
   \bigl(\mathbf H^{-q_1}(X^{(1)},\mathscr P_\bullet^{(1)}),
         \mathbf H^{-q_2}(X^{(2)},\mathscr P_\bullet^{(2)})\bigr).
\]
Por definición, el complejo
$\mathrm H_n^{\mathrm I}(L_{\bullet\bullet}^{(i)})$ tiene por término de grado
$k$ el módulo de homología
$\mathrm H_n(C^\bullet(\mathfrak U^{(i)},\mathscr P_k^{(i)}))$, es decir,
por definición, el módulo de cohomología
$\mathbf H^{-n}(\mathfrak U^{(i)},\mathscr P_k^{(i)})$; se sabe por (1.4.1)
que este módulo es canónicamente isomorfo a
$\mathbf H^{-n}(X^{(i)},\mathscr P_k^{(i)})$; por tanto,
\[
  {}^{(a')}\mathscr E_{pq}^2
  =\bigoplus_{q_1+q_2=q}
   \mathbf{Tor}_p^A
   \bigl(\mathbf H^{-q_1}(X^{(1)},\mathscr P_\bullet^{(1)}),
         \mathbf H^{-q_2}(X^{(2)},\mathscr P_\bullet^{(2)})\bigr).
\]
\end{env}

\begin{env}[6.6.4]
\label{III.6.6.4}
\emph{Sucesiones espectrales $(b)$ y $(b')$.} Por definición,
$\mathbf{Tor}_q^{A,\mathrm{II}}
(L_{\bullet\bullet}^{(1)},L_{\bullet\bullet}^{(2)})$ es un bicomplejo cuya
componente de bigrado $(h,k)$ es el $A$-módulo
\[
  \mathbf{Tor}_q^A
  \bigl(C^{-h}(\mathfrak U^{(1)},\mathscr P_\bullet^{(1)}),
        C^{-k}(\mathfrak U^{(2)},\mathscr P_\bullet^{(2)})\bigr).
\]
Sea $\Phi^{(i)}$ el conjunto de índices de $\mathfrak U^{(i)}$; por
definición, el complejo de módulos
$C^r(\mathfrak U^{(i)},\mathscr P_\bullet^{(i)})$ ($r\geqslant0$) es suma
directa de los complejos
$\Gamma(U_\rho^{(i)},\mathscr P_\bullet^{(i)})$, donde $U_\rho^{(i)}$ es la
intersección de los $U_\xi^{(i)}$ para $\xi\in\rho$, y $\rho$ recorre
$\mathfrak P(\Phi^{(i)})$; por tanto, el $A$-módulo
\[
  \mathbf{Tor}_q^A
  \bigl(C^{-h}(\mathfrak U^{(1)},\mathscr P_\bullet^{(1)}),
        C^{-k}(\mathfrak U^{(2)},\mathscr P_\bullet^{(2)})\bigr)
\]
es suma directa de los $A$-módulos
\[
  \mathbf{Tor}_q^A
  \bigl(\Gamma(U_\sigma^{(1)},\mathscr P_\bullet^{(1)}),
        \Gamma(U_\tau^{(2)},\mathscr P_\bullet^{(2)})\bigr),
\]
donde $\sigma$ (respectivamente, $\tau$) recorre los elementos de
$\mathfrak P(\Phi^{(1)})$ (respectivamente, $\mathfrak P(\Phi^{(2)})$) tales
que $\operatorname{Card}(\sigma)=-(h+1)$ (respectivamente,
$\operatorname{Card}(\tau)=-(k+1)$). Como $X^{(1)}$ y $X^{(2)}$ son
esquemas, los $U_\rho^{(i)}$ son afines y, por tanto, se tiene (6.4.1.1)
\[
  \mathbf{Tor}_q^A
  \bigl(\Gamma(U_\sigma^{(1)},\mathscr P_\bullet^{(1)}),
        \Gamma(U_\tau^{(2)},\mathscr P_\bullet^{(2)})\bigr)
  =\Gamma\bigl(U_\sigma^{(1)}\times_S U_\tau^{(2)},
    \mathscr{T}\!or_q^S
    (\mathscr P_\bullet^{(1)},\mathscr P_\bullet^{(2)})\bigr).
\]
\oldpage[III]{155}
Se ve, pues, que ${}^{(b)}\mathscr E_{pq}^2$ es el $(-p)$-ésimo módulo de
cohomología del complejo
\[
  L^\bullet(\Phi^{(1)},\Phi^{(2)};\mathscr S)
\]
de cocadenas bialternadas sobre $\Phi^{(1)}$ y $\Phi^{(2)}$ con valores en el
sistema de coeficientes
\[
  \mathscr S:(\sigma,\tau)\leadsto
  \Gamma\bigl(U_\sigma^{(1)}\times_S U_\tau^{(2)},
    \mathscr{T}\!or_q^S
    (\mathscr P_\bullet^{(1)},\mathscr P_\bullet^{(2)})\bigr)
\]
($0_{\mathrm I}$, 11.8.4). Se sabe entonces ($0_{\mathrm I}$, 11.8.5 y
11.8.6) que la cohomología de este complejo es la misma que la del complejo
$C^\bullet(\Phi^{(1)},\Phi^{(2)};\mathscr S)$ de todas las cocadenas sobre
$\Phi^{(1)}$ y $\Phi^{(2)}$ con valores en $\mathscr S$, y también la del
complejo $P^\bullet(\Phi^{(1)},\Phi^{(2)};\mathscr S)$, cuyos elementos son
las combinaciones lineales de los
\[
  \lambda(\sigma,\tau)\in
  \Gamma\bigl(U_\sigma^{(1)}\times_S U_\tau^{(2)},
    \mathscr{T}\!or_q^S
    (\mathscr P_\bullet^{(1)},\mathscr P_\bullet^{(2)})\bigr),
\]
donde $\sigma=(\alpha_0,\ldots,\alpha_h)$ y
$\tau=(\beta_0,\ldots,\beta_h)$ son sucesiones con el mismo número de
elementos. Pero entonces se tiene
\[
  U_\sigma^{(1)}\times_S U_\tau^{(2)}
  =(U_{\alpha_0}^{(1)}\times_S U_{\beta_0}^{(2)})\cap\cdots\cap
    (U_{\alpha_h}^{(1)}\times_S U_{\beta_h}^{(2)})
  \qquad\text{(I, 3.2.7)}.
\]
Si se denota por $\mathfrak U$ el cubrimiento de
$Z=X^{(1)}\times_S X^{(2)}$ por los abiertos afines
$U_\alpha^{(1)}\times_S U_\beta^{(2)}$, se ve finalmente, teniendo en cuenta
que $X^{(1)}\times_S X^{(2)}$ es un esquema, que se tiene, en virtud de
(1.3.1),
\[
  {}^{(b)}\mathscr E_{pq}^2
  =\mathbf H^{-p}\bigl(X^{(1)}\times_S X^{(2)},
    \mathscr{T}\!or_q^S
    (\mathscr P_\bullet^{(1)},\mathscr P_\bullet^{(2)})\bigr).
\]

En segundo lugar,
$\mathbf{Tor}_q^A(L_{\bullet\bullet}^{(1)},L_{\bullet\bullet}^{(2)})$ es un
bicomplejo cuya componente de bigrado $(h,k)$ es la suma directa de los
$A$-módulos
\[
  \mathbf{Tor}_q^A
  \bigl(C^{-h_1}(\mathfrak U^{(1)},\mathscr P_{k_1}^{(1)}),
        C^{-h_2}(\mathfrak U^{(2)},\mathscr P_{k_2}^{(2)})\bigr)
\]
tales que $h_1+h_2=h$ y $k_1+k_2=k$; explicitando los módulos
$C^r(\mathfrak U^{(i)},\mathscr P_s^{(i)})$ como antes, se ve nuevamente que
esta componente es suma directa de los $A$-módulos
\[
  \Gamma\bigl(U_\sigma^{(1)}\times_S U_\tau^{(2)},
    \mathscr{T}\!or_q^S
    (\mathscr P_{k_1}^{(1)},\mathscr P_{k_2}^{(2)})\bigr),
\]
donde $k_1+k_2=k$, y $\sigma$ (respectivamente, $\tau$) recorre los elementos
de $\mathfrak P(\Phi^{(1)})$ (respectivamente,
$\mathfrak P(\Phi^{(2)})$) tales que
$\operatorname{Card}(\sigma)+\operatorname{Card}(\tau)=-h-2$. El término
${}^{(b')}\mathscr E_{pq}^2$ que estamos calculando es el $(-p)$-ésimo módulo
de cohomología de un bicomplejo $N^{\bullet\bullet}=(N^{hk})$, donde el
complejo simple $N^{\bullet k}$ es el complejo de cocadenas bialternadas sobre
$\Phi^{(1)}$ y $\Phi^{(2)}$, con valores en el sistema de coeficientes
\[
  \mathscr S_k:(\sigma,\tau)\leadsto
  \Gamma\left(U_\sigma^{(1)}\times_S U_\tau^{(2)},
    \bigoplus_{k_1+k_2=k}
    \mathscr{T}\!or_q^S
    (\mathscr P_{-k_1}^{(1)},\mathscr P_{-k_2}^{(2)})\right),
\]
y estos sistemas de coeficientes forman un complejo $\mathscr S^\bullet$ cuyo
diferencial proviene del diferencial del complejo simple asociado al bicomplejo
$\mathscr{T}\!or_q^S
(\mathscr P_\bullet^{(1)},\mathscr P_\bullet^{(2)})$. Se sabe que la
cohomología de $N^{\bullet\bullet}$ es la misma que la del bicomplejo
$C^\bullet(\Phi^{(1)},\Phi^{(2)};\mathscr S^\bullet)$
($0_{\mathrm I}$, 11.8.9), y también la misma que la del bicomplejo
$P^\bullet(\Phi^{(1)},\Phi^{(2)};\mathscr S^\bullet)$, donde los elementos de
bigrado $(h,k)$ son las combinaciones lineales de los
\[
  \lambda(\sigma,\tau)\in
  \Gamma\left(U_\sigma^{(1)}\times_S U_\tau^{(2)},
    \bigoplus_{k_1+k_2=k}
    \mathscr{T}\!or_q^S
    (\mathscr P_{-k_1}^{(1)},\mathscr P_{-k_2}^{(2)})\right),
\]
donde $\sigma=(\alpha_0,\ldots,\alpha_h)$ y
$\tau=(\beta_0,\ldots,\beta_h)$ son sucesiones con el mismo número de
elementos ($0_{\mathrm I}$, 11.8.10). Se ve entonces, como antes, que
${}^{(b')}\mathscr E_{pq}^2$ es el $(-p)$-ésimo módulo de cohomología del
bicomplejo $C^\bullet(\mathfrak U,\mathscr Q^\bullet)$, donde
$\mathscr Q^\bullet$ es el complejo simple asociado al bicomplejo
$\mathscr{T}\!or_q^S
(\mathscr P_\bullet^{(1)},\mathscr P_\bullet^{(2)})$ de
$\mathscr O_Z$-módulos. Con las convenciones adoptadas en (6.6.1), se tiene,
por tanto,
\[
  {}^{(b')}\mathscr E_{pq}^2
  =\mathbf H^{-p}\bigl(X^{(1)}\times_S X^{(2)},
    \mathscr{T}\!or_q^S
    (\mathscr P_\bullet^{(1)},\mathscr P_\bullet^{(2)})\bigr).
\]
\end{env}

\oldpage[III]{156}
\begin{env}[6.6.5]
\label{III.6.6.5}
\emph{Sucesiones espectrales $(c)$ y $(d)$.} Por definición, el complejo
$\mathrm H_n^{\mathrm{II}}(L_{\bullet\bullet}^{(i)})$ tiene por término de
grado $h$ el $A$-módulo
$C^{-h}(\mathfrak U^{(i)},\mathscr H_n(\mathscr P_\bullet^{(i)}))$, en virtud
de la exactitud del funtor $C^{-h}$. Por tanto, por definición del hipertor de
dos módulos relativo a dos cubrimientos (6.6.2),
\[
  {}^{(c)}\mathscr E_{pq}^2
  =\bigoplus_{q_1+q_2=q}
   \mathbf{Tor}_p^S
   \bigl(\mathfrak U^{(1)},\mathfrak U^{(2)};
         \mathscr H_{q_1}(\mathscr P_\bullet^{(1)}),
         \mathscr H_{q_2}(\mathscr P_\bullet^{(2)})\bigr).
\]
Por último, por definición,
$\mathbf{Tor}_q^{A,\mathrm I}
(L_{\bullet\bullet}^{(1)},L_{\bullet\bullet}^{(2)})$ es un bicomplejo cuya
componente de bigrado $(h,k)$ es el $A$-módulo
$\mathbf{Tor}_q^S
(\mathfrak U^{(1)},\mathfrak U^{(2)};
 \mathscr P_h^{(1)},\mathscr P_k^{(2)})$. Por tanto,
\[
  {}^{(d)}\mathscr E_{pq}^2
  =\mathrm H_p\bigl(
    \mathbf{Tor}_q^S
    (\mathfrak U^{(1)},\mathfrak U^{(2)};
     \mathscr P_\bullet^{(1)},\mathscr P_\bullet^{(2)})\bigr).
\]
\end{env}

\begin{env}[6.6.6]
\label{III.6.6.6}
La teoría de la hiperhomología de los funtores de bicomplejos
($0_{\mathrm I}$, 11.7.3) muestra, como en (6.3.4), que, para toda resolución
\emph{plana} de Cartan--Eilenberg $M_{\bullet\bullet}^{(i)}$ de
$L_{\bullet\bullet}^{(i)}$ (en la categoría de los complejos de módulos
acotados inferiormente) ($i=1,2$), se tienen isomorfismos canónicos de
bi-$\partial$-funtores
\begin{env}[6.6.6.1]
\label{III.6.6.6.1}
\[
\begin{aligned}
  \mathbf{Tor}_\bullet^S
  (\mathfrak U^{(1)},\mathfrak U^{(2)};
   \mathscr P_\bullet^{(1)},\mathscr P_\bullet^{(2)})
  &\xrightarrow{\sim}
   \mathrm H_\bullet
   (M_{\bullet\bullet}^{(1)}\otimes_A M_{\bullet\bullet}^{(2)})\\
  &\xrightarrow{\sim}
   \mathrm H_\bullet
   (M_{\bullet\bullet}^{(1)}\otimes_A L_{\bullet\bullet}^{(2)})
   \xrightarrow{\sim}
   \mathrm H_\bullet
   (L_{\bullet\bullet}^{(1)}\otimes_A M_{\bullet\bullet}^{(2)}).
\end{aligned}
\tag{6.6.6.1}
\]
\end{env}
\end{env}

\begin{env}[6.6.7]
\label{III.6.6.7}
Vamos a mostrar ahora que el hipertor global definido en (6.6.2), así como
las seis sucesiones espectrales correspondientes, no dependen de los
cubrimientos abiertos afines finitos $\mathfrak U^{(i)}$ que se utilizaron
para definirlos (salvo isomorfismos canónicos). Bastará para ello mostrar que,
si $\mathfrak V^{(i)}$ son otros dos cubrimientos de la misma naturaleza,
tales que $\mathfrak V^{(i)}$ sea \emph{más fino} que $\mathfrak U^{(i)}$
para $i=1,2$, entonces se tienen \emph{isomorfismos canónicos} de funtores
espectrales
\begin{env}[6.6.7, $t$]
\label{III.6.6.7.t}
\[
  {}^{(t)}\mathscr E
  (\mathfrak U^{(1)},\mathfrak U^{(2)};
   \mathscr P_\bullet^{(1)},\mathscr P_\bullet^{(2)})
  \xrightarrow{\sim}
  {}^{(t)}\mathscr E
  (\mathfrak V^{(1)},\mathfrak V^{(2)};
   \mathscr P_\bullet^{(1)},\mathscr P_\bullet^{(2)})
  \tag{6.6.7, $t$}\label{III.6.6.7.t.eq}
\]
\end{env}
donde $t$ se sustituye por $a$, $b$, $a'$, $b'$, $c$ o $d$.

Ahora bien, para $i=1,2$ se tienen homomorfismos de bicomplejos
\[
  C^\bullet(\mathfrak U^{(i)},\mathscr P_\bullet^{(i)})
  \longrightarrow
  C^\bullet(\mathfrak V^{(i)},\mathscr P_\bullet^{(i)})
\]
bien definidos salvo homotopía (G, II, 5.7.1); de ello resultan ya
homomorfismos (6.6.7, $t$) definidos canónicamente y compatibles con los
operadores de borde en los términos límite ($0_{\mathrm I}$, 11.3.2). Además,
el cálculo de los términos $E_2$ de las sucesiones espectrales $(a)$, $(b)$,
$(a')$, $(b')$ muestra que, para estas sucesiones espectrales, el homomorfismo
(6.6.7, $t$) es un isomorfismo sobre los términos $E_2$; como estas sucesiones
espectrales son birregulares, se ve que (6.6.7, $t$) es un isomorfismo para
estos funtores espectrales y, por consiguiente, induce un \emph{isomorfismo de
bi-$\partial$-funtores} en su término límite común
($0_{\mathrm I}$, 11.1.5).

En particular, para $\mathscr O_{X^{(i)}}$-módulos cuasicoherentes
$\mathscr F^{(i)}$ ($i=1,2$), el homomorfismo canónico
\[
  \mathbf{Tor}_\bullet^S
  (\mathfrak U^{(1)},\mathfrak U^{(2)};
   \mathscr F^{(1)},\mathscr F^{(2)})
  \longrightarrow
  \mathbf{Tor}_\bullet^S
  (\mathfrak V^{(1)},\mathfrak V^{(2)};
   \mathscr F^{(1)},\mathscr F^{(2)})
\]
es biyectivo; visto el cálculo de (6.6.5), se ve que (6.6.7, $t$) es también
un isomorfismo de los términos $E_2$ para $t=c$ y $t=d$. Se concluye como
antes que (6.6.7, $t$) es también un isomorfismo de sucesiones espectrales
para $t=c$ y $t=d$.
\oldpage[III]{157}
Se puede considerar que los isomorfismos (6.6.7, $t$) definen sistemas
inductivos de funtores espectrales sobre el conjunto filtrante de pares
$(\mathfrak U^{(1)},\mathfrak U^{(2)})$ de cubrimientos abiertos afines
finitos de $X^{(1)}$ y $X^{(2)}$. Denotaremos por
\[
  {}^{(t)}\mathscr E
  (X^{(1)},X^{(2)};
   \mathscr P_\bullet^{(1)},\mathscr P_\bullet^{(2)})
  \quad\text{o}\quad
  {}^{(t)}\mathscr E^S
  (X^{(1)},X^{(2)};
   \mathscr P_\bullet^{(1)},\mathscr P_\bullet^{(2)})
\]
el límite inductivo de este sistema, y por
$\mathbf{Tor}_\bullet^S
(X^{(1)},X^{(2)};\mathscr P_\bullet^{(1)},\mathscr P_\bullet^{(2)})$
el término límite de este funtor espectral, al que llamaremos el
\emph{hipertor global} de los dos complejos $\mathscr P_\bullet^{(1)}$ y
$\mathscr P_\bullet^{(2)}$; si $\mathscr P_\bullet^{(1)}$ y
$\mathscr P_\bullet^{(2)}$ están reducidos a sus términos de grado 0,
$\mathscr F^{(1)}$ y $\mathscr F^{(2)}$, escribiremos
\[
  \mathbf{Tor}_\bullet^S
  (X^{(1)},X^{(2)};\mathscr F^{(1)},\mathscr F^{(2)}),
\]
y, conforme a las convenciones generales,
\[
  \mathbf{Tor}_q^S
  (X^{(1)},X^{(2)};\mathscr P_\bullet^{(1)},\mathscr P_\bullet^{(2)})
\]
será, por tanto, el bicomplejo cuyas componentes son los módulos
\[
  \mathbf{Tor}_q^S
  (X^{(1)},X^{(2)};\mathscr P_h^{(1)},\mathscr P_k^{(2)}).
\]
\end{env}

\begin{env}[6.6.8]
\label{III.6.6.8}
Bajo las hipótesis de (6.6.1), consideremos ahora dos $S$-morfismos
$f_i:X^{(i)}\to Y^{(i)}$, donde
$Y^{(i)}=\operatorname{Spec}(B_i)$ es un $S$-esquema \emph{afín}, siendo
$B_i$ una $A$-álgebra ($i=1,2$); esto define, pues, un $A$-homomorfismo
$B_i\to\Gamma(X^{(i)},\mathscr O_{X^{(i)}})$ (I, 2.2.4) y, por consiguiente,
cada uno de los $L_{\bullet\bullet}^{(i)}$ definidos en (6.6.2) es un
bicomplejo de $B_i$-módulos; se concluye que
$L_{\bullet\bullet}^{(1)}\otimes_A L_{\bullet\bullet}^{(2)}$ es un
complejo cuádruple de $(B_1\otimes_A B_2)$-módulos, y sus seis funtores
espectrales de hiperhomología pueden considerarse, por tanto, con valores en
la categoría de las sucesiones espectrales de
$(B_1\otimes_A B_2)$-módulos. Si se pone
$Y=Y^{(1)}\times_S Y^{(2)}=\operatorname{Spec}(B_1\otimes_A B_2)$, se pueden
considerar los $\mathscr O_Y$-módulos cuasicoherentes asociados a estos
módulos (I, 1.3.4); denotaremos por
${}^{(t)}\mathscr E
(f_1,f_2;\mathscr P_\bullet^{(1)},\mathscr P_\bullet^{(2)})$
(para $t=a$, $b$, $a'$, $b'$, $c$ o $d$) las seis sucesiones espectrales
\[
  \bigl({}^{(t)}\mathscr E
  (X^{(1)},X^{(2)};
   \mathscr P_\bullet^{(1)},\mathscr P_\bullet^{(2)})\bigr)^{\sim}
\]
de $\mathscr O_Y$-módulos, y por
$\mathfrak{Tor}_\bullet^S
(f_1,f_2;\mathscr P_\bullet^{(1)},\mathscr P_\bullet^{(2)})$ su término
límite común
\[
  \bigl(\mathbf{Tor}_\bullet^S
  (X^{(1)},X^{(2)};
   \mathscr P_\bullet^{(1)},\mathscr P_\bullet^{(2)})\bigr)^{\sim}.
\]
Lo denotaremos por
$\mathfrak{Tor}_\bullet^S(f_1,f_2;\mathscr F^{(1)},\mathscr F^{(2)})$
cuando $\mathscr P_\bullet^{(i)}$ esté reducido a su término de grado 0,
$\mathscr F^{(i)}$ ($i=1,2$).
\end{env}
